\documentclass[11pt]{article}
\newcommand{\SheetCode}{Theory 05}
\usepackage[margin=25mm]{geometry}
\usepackage[T1]{fontenc}
\usepackage{lmodern}
\DeclareUnicodeCharacter{2605}{$\star$}
\usepackage{microtype}
\usepackage{amsmath,amssymb,mathtools,bm}
\usepackage{enumitem}
\usepackage{xcolor}
\usepackage{fancyhdr}
\usepackage{hyperref}
\usepackage[most]{tcolorbox}
\definecolor{agiBlue}{HTML}{173B57}
\definecolor{agiGold}{HTML}{B7791F}
\definecolor{agiLight}{HTML}{EFF6FA}
\definecolor{agiGreen}{HTML}{184D3B}
\hypersetup{colorlinks=true,linkcolor=agiBlue,urlcolor=agiBlue}
\setlength{\parindent}{0pt}
\setlength{\parskip}{0.55em}
\setlength{\headheight}{14pt}
\setlist{nosep,leftmargin=*}
\pagestyle{fancy}
\fancyhf{}
\fancyhead[L]{\small Part of the \href{https://github.com/mmikhasenko/GIModel.jl}{Agentic Gottried--Isgur} project}
\fancyhead[R]{\SheetCode}
\fancyfoot[C]{\thepage}
\newcommand{\dd}{\,\mathrm d}
\newcommand{\ket}[1]{\lvert #1\rangle}
\newcommand{\bra}[1]{\langle #1\rvert}
\newcommand{\braket}[2]{\langle #1\mid #2\rangle}
\newcommand{\expect}[1]{\left\langle #1\right\rangle}
\newcommand{\op}[1]{\widehat{#1}}
\newcommand{\GeV}{\,\mathrm{GeV}}
\newcommand{\R}{\mathbb{R}}
\newtcolorbox{goals}{colback=agiLight,colframe=agiBlue,title=Learning outcomes}
\newtcolorbox{reading}{colback=white,colframe=agiGold,title=Book reference,
  after skip=5mm}
\newtcolorbox{paperbridge}{colback=agiGreen!7,colframe=agiGreen,
  colbacktitle=agiGreen,coltitle=white,title=Connection to the GI paper}
\newtcolorbox{problem}[1]{colback=white,colframe=agiBlue,title={#1},
  before skip=3.5mm,after skip=1mm}
\newtcolorbox{solution}{colback=black!2,colframe=black!45,title=Solution}
\newtcolorbox{quiz}{colback=white,colframe=agiGold,title=Post-solution concept check}
\newtcolorbox{checkpoint}{colback=white,colframe=agiGold,title=Interpretation checkpoint}
\newcommand{\sheettitle}[3]{%
  \begin{center}
  {\Large\bfseries #1}\par
  \vspace{0.2em}{\color{agiBlue}#2}\par
  \vspace{0.4em}\small #3
  \end{center}\vspace{0.5em}}
\newcommand{\difficulty}[1]{\textbf{Difficulty:} #1}
\newcommand{\timebox}[1]{\textbf{Suggested time:} #1}
\newcommand{\answer}{\par\smallskip\color{black!50}\textbf{Answer:} }

\begin{document}
\sheettitle{Sheet 5: Making singular interactions well-defined}{Gaussian smearing and Hermitian momentum sandwiches}{Prerequisite: Sheets 1--4. \difficulty{★★--★★★} \quad \timebox{3--4 hours}}

\begin{goals}
Normalize the GI Gaussian; derive its convolution rules; understand the
mass-dependent resolution scale; prove Hermiticity and positivity of momentum
sandwiches; distinguish an operator product from a pointwise multiplier.
\end{goals}
\begin{reading}
Selected reading in Landau--Lifshitz, \emph{Quantum Mechanics: Non-Relativistic
Theory}: Ch. I, ``The Basic Concepts of Quantum Mechanics,'' and Ch. II,
``Energy and Momentum,'' for operators and momentum representation. This is an
optional reference selection; the sheet introduces the needed constructions.
\end{reading}


Define
$\rho_\sigma(\bm r)=\sigma^3\pi^{-3/2}e^{-\sigma^2r^2}$ and
$\widetilde f(\bm r)=\int\rho_\sigma(\bm r-\bm r')f(\bm r')\dd^3r'$.
The phrase \emph{momentum sandwich} below is shorthand for the ordered operator
product $B(p)V(r)B(p)$: the same momentum-dependent factor is placed on both
sides of a position-dependent interaction.

\begin{problem}{Problem 1 ★★ — Normalize and measure the smearing kernel}
Prove $\int\rho_\sigma\dd^3r=1$ and compute $\expect{r^2}_\rho$.
\end{problem}
\begin{solution}
The three Cartesian Gaussian integrals give
$(\sqrt\pi/\sigma)^3$, cancelling the prefactor. Each coordinate has variance
$1/(2\sigma^2)$, hence $\expect{r^2}=3/(2\sigma^2)$. Thus larger $\sigma$
means a narrower spatial resolution.
\end{solution}
\begin{quiz}
Does $\sigma\to\infty$ represent stronger or weaker smearing? \answer weaker;
$\rho_\sigma$ approaches $\delta^3(\bm r)$.
\end{quiz}

\begin{problem}{Problem 2 ★★ — Smear a Gaussian}
For $f(r)=e^{-\gamma^2r^2}$, compute $\widetilde f(r)$ and identify the combined
width.
\end{problem}
\begin{solution}
Completing the square in the convolution gives
\[
\widetilde f(r)=\left(\frac{\sigma^2}{\sigma^2+\gamma^2}\right)^{3/2}
\exp\!\left[-\frac{\sigma^2\gamma^2}{\sigma^2+\gamma^2}r^2\right].
\]
If $\tau^{-2}=\sigma^{-2}+\gamma^{-2}$, the exponential is
$e^{-\tau^2r^2}$.
\end{solution}
\begin{quiz}
Why are inverse squared widths added? \answer convolution adds coordinate-space
variances, and a Gaussian variance scales as the inverse squared width.
\end{quiz}

\begin{problem}{Problem 3 ★★ — Check the GI mass-dependent width}
GI uses
\[
\sigma^2=\sigma_0^2\left[\frac12+\frac12
\left(\frac{4m_1m_2}{(m_1+m_2)^2}\right)^4\right]
+s^2\left(\frac{2m_1m_2}{m_1+m_2}\right)^2.
\]
Compute the equal-mass form and the $m_1\gg m_2$ limit. Interpret both.
\end{problem}
\begin{solution}
For $m_1=m_2=m$, the bracket equals one and the reduced combination equals $m$,
so $\sigma^2=\sigma_0^2+s^2m^2$. For $m_1\gg m_2$, the fourth-power term
vanishes and $2m_1m_2/(m_1+m_2)\to2m_2$, giving
$\sigma^2\to\sigma_0^2/2+4s^2m_2^2$. The light constituent controls the
resolution in a heavy--light system.
\end{solution}
\begin{quiz}
Why is a single universal smearing radius inadequate across all flavor sectors?
\answer relativistic localization is mass-dependent; light and heavy pairs
resolve short distances differently.
\end{quiz}

\begin{problem}{Problem 4 ★★ — Prove a momentum sandwich is Hermitian}
Let $B=B(p)$ be a real function of the self-adjoint $p^2$, and $V=V(r)$ a real
multiplication operator. Prove $BVB$ is Hermitian and show
$\bra\psi BVB\ket\psi=\bra{B\psi}V\ket{B\psi}$.
\end{problem}
\begin{solution}
$B^\dagger=B$ and $V^\dagger=V$, so
$(BVB)^\dagger=BVB$. Associativity gives the expectation identity. If $V$ is
positive, the right-hand side is nonnegative, so the sandwich also preserves
positivity.
\end{solution}
\begin{quiz}
Would $BV$ generally be Hermitian? \answer no; $B(p)$ and $V(r)$ do not commute.
The two-sided ordering is physical structure, not cosmetic symmetrization.
\end{quiz}

\begin{problem}{Problem 5 ★★★ — Analyze the GI relativization factor}
For $E_i(p)=\sqrt{p^2+m_i^2}$, consider
$B_i(p)=[m_1m_2/(E_1E_2)]^{1/2+\epsilon_i}$ and
$A(p)=[1+p^2/(E_1E_2)]^{1/2}$. Compute their $p\to0$ and $p\to\infty$ limits
for positive masses and explain their different roles.
\end{problem}
\begin{solution}
At $p=0$, $B_i=A=1$. At large $p$, $E_1E_2\sim p^2$, so
$A\to\sqrt2$, whereas
$B_i\sim(m_1m_2/p^2)^{1/2+\epsilon_i}$. $A$ dresses the central vector
interaction without suppressing high momentum; $B_i$ relativizes inverse-mass
spin operators and softens their high-momentum action. The exponent is attached
to each side of the Hermitian sandwich.
\end{solution}
\begin{quiz}
Why can one not reproduce $BVB$ by multiplying $V(r)$ by an average number?
\answer $B$ reshapes the state in momentum space before and after $V$; the
result depends on off-diagonal matrix elements and on the state.
\end{quiz}

\begin{paperbridge}
GI Eq. (10), repeated as Appendix-A Eq. (A7), defines the normalized Gaussian;
(A8)--(A9) define convolution and its mass-dependent width. Eqs. (A12)--(A14)
give the closed smeared potentials. The prescriptions immediately before
(A15), and then (A15)--(A16), place the momentum-dependent factors on both sides
of each radial interaction.
\end{paperbridge}
\end{document}
